% =====================================================================
%  1. Introducción
% =====================================================================
\section{Introducción}
\label{sec:introduccion}

Este documento reporta el estado de una investigación en curso sobre
\textbf{EduCurator}, un sistema de curación asistida de documentación de cursos
universitarios. La investigación tiene una particularidad que conviene declarar
desde el principio: el artefacto tecnológico \emph{ya existe}. Hay un producto
mínimo viable (MVP) funcional que implementa la arquitectura propuesta y que puede
ejecutar el flujo completo de curación sobre un corpus real. Lo que no existe
todavía es la evidencia empírica sobre su desempeño.

Esa asimetría define la naturaleza del documento. No se trata de un informe de
resultados ni de la memoria de un desarrollo de software, sino de un
\textbf{protocolo de investigación acompañado de un reporte de avance}: establece
qué se quiere averiguar sobre el artefacto, con qué diseño experimental, con qué
instrumentos y bajo qué criterios de análisis; y, en paralelo, delimita con
precisión qué está construido y qué permanece abierto.

\subsection{Por qué el artefacto no es el resultado}

Un error frecuente en los proyectos de ingeniería con componente investigativo es
presentar el sistema construido como si fuera el hallazgo de la investigación. Una
funcionalidad implementada demuestra \emph{capacidad técnica}: que el sistema puede
ejecutar una operación. Una evaluación demuestra \emph{desempeño}: con qué
exactitud, bajo qué condiciones y frente a qué alternativas. Son afirmaciones de
distinto tipo y requieren distinta evidencia.

La metodología de \emph{Design Science Research} permite sostener esta distinción
sin contradicción metodológica. \textcite{peffers2007design} admiten explícitamente
puntos de entrada distintos al ciclo de investigación, entre ellos uno
\emph{centrado en el diseño y desarrollo}, en el que el artefacto precede a la
formalización del problema y de la evaluación. \textcite{hevner2004design} exigen,
a cambio, que la contribución se demuestre mediante métodos de evaluación
rigurosos: el artefacto por sí solo no constituye conocimiento. Este documento se
sitúa exactamente en ese punto: el diseño y el desarrollo están hechos; la
demostración y la evaluación son el trabajo que falta.

\subsection{Qué pregunta abre el problema}

La curación de material docente --- comparar versiones, detectar duplicados,
identificar contradicciones, reconocer lo que quedó desactualizado --- es una tarea
que crece con el volumen del corpus y que rara vez se completa de forma sistemática.
La expectativa razonable sería que la inteligencia artificial generativa la aliviara.
La evidencia sistemática disponible advierte contra esa expectativa: en una revisión
PRISMA de 30 estudios empíricos, \textcite{su2026efficiency} encuentran que solo el
20\,\% reporta reducción de la carga docente, el 16,7\,\% reporta aumento y el
63,3\,\% reporta un efecto mixto. Su conclusión es que la IA generativa produce una
transformación estructural más que una reducción:
\textcquote[p.~11]{su2026efficiency}{the key change brought about by GenAI
is not whether work becomes \enquote{less}, but how workload is redistributed}.

El mecanismo que describen es directamente pertinente aquí: la IA ahorra tiempo en
tareas procedimentales, pero añade trabajo de revisión de calidad de sus salidas,
verificación de hechos y juicio sobre lo que propone. Es decir, un asistente de
curación que proponga cambios sin evidencia verificable traslada al docente el costo
de comprobarlos y puede terminar siendo una fuente neta de trabajo. Ese es el riesgo
de diseño que EduCurator intenta atacar y, a la vez, la razón por la que su
evaluación no puede limitarse a medir tiempo.

\subsection{Función y organización del documento}

El documento está organizado en siete partes:

\begin{itemize}
  \item La \textbf{Parte~I} sitúa el problema, lo delimita y lo justifica.
  \item La \textbf{Parte~II} revisa el conocimiento existente en las ocho líneas
        bibliográficas del proyecto y formula la brecha de investigación.
  \item La \textbf{Parte~III} enuncia la pregunta, los objetivos --- separando los
        de investigación de los de desarrollo ---, las proposiciones, las variables
        operacionalizadas y el modelo conceptual que las vincula.
  \item La \textbf{Parte~IV} describe el artefacto: qué propone, qué implementa y
        qué no se sabe aún sobre él.
  \item La \textbf{Parte~V} define la metodología de evaluación: benchmark,
        configuraciones de comparación, ablaciones, instrumentos, métricas y plan de
        análisis.
  \item La \textbf{Parte~VI} reporta el estado de avance, la evidencia técnica
        disponible y lo pendiente.
  \item La \textbf{Parte~VII} enuncia resultados esperados, contribuciones
        esperadas y el cronograma de validación.
\end{itemize}

\subsection{Criterio de lectura}
\label{sec:criterio-lectura}

Para mantener separadas afirmaciones de distinto estatus epistémico, cada
enunciado sustantivo del documento corresponde a una de cuatro categorías:

\begin{enumerate}[label=\textbf{(\Alph*)}]
  \item \textbf{Lo que la literatura ya demuestra}, con su fuente.
  \item \textbf{Lo que el MVP actualmente implementa}, verificable en el
        repositorio.
  \item \textbf{Lo que la investigación propone evaluar}, con su diseño.
  \item \textbf{Lo que todavía no se conoce} y constituye el objeto de la
        investigación.
\end{enumerate}

\begin{advertencia}
Ningún enunciado de este documento afirma que EduCurator mejore la curación,
reduzca el tiempo docente, elimine las alucinaciones o supere a alternativa alguna.
Esas serían afirmaciones de la categoría (A) sin fuente que las respalde, o
afirmaciones de la categoría (C) presentadas como si ya fueran (A). Donde el texto
se refiere al comportamiento del sistema, lo hace en términos de lo que
\emph{implementa} o de lo que \emph{se propone medir}.
\end{advertencia}
